\documentclass[10pt,a4paper]{amsart}
\usepackage{amsmath,amssymb,amsthm,graphicx}
\usepackage{verbatim}
\newenvironment{reference}{\comment}{\endcomment}
%\newenvironment{reference}{}{}
\newenvironment{slogan}{\comment}{\endcomment}
\newenvironment{history}{\comment}{\endcomment}

% For commutative diagrams we use Xy-pic
\usepackage[all]{xy}

% We use 2cell for 2-commutative diagrams.
\xyoption{2cell}
% Theorem environments.
%
\theoremstyle{plain}
\newtheorem{theorem}[subsection]{Theorem}
\newtheorem{proposition}[subsection]{Proposition}
\newtheorem{lemma}[subsection]{Lemma}

\theoremstyle{definition}
\newtheorem{definition}[subsection]{Definition}
\newtheorem{example}[subsection]{Example}
\newtheorem{exercise}[subsection]{Exercise}
\newtheorem{situation}[subsection]{Situation}

\theoremstyle{remark}
\newtheorem{remark}[subsection]{Remark}
\newtheorem{remarks}[subsection]{Remarks}

\numberwithin{equation}{subsection}

% Macros
%
\def\lim{\mathop{\mathrm{lim}}\nolimits}
\def\colim{\mathop{\mathrm{colim}}\nolimits}
\def\Spec{\mathop{\mathrm{Spec}}}
\def\Hom{\mathop{\mathrm{Hom}}\nolimits}
\def\Ext{\mathop{\mathrm{Ext}}\nolimits}
\def\SheafHom{\mathop{\mathcal{H}\!\mathit{om}}\nolimits}
\def\SheafExt{\mathop{\mathcal{E}\!\mathit{xt}}\nolimits}
\def\Sch{\mathit{Sch}}
\def\Mor{\mathop{\mathrm{Mor}}\nolimits}
\def\Ob{\mathop{\mathrm{Ob}}\nolimits}
\def\Sh{\mathop{\mathit{Sh}}\nolimits}
\def\NL{\mathop{N\!L}\nolimits}
\def\CH{\mathop{\mathrm{CH}}\nolimits}
\def\proetale{{pro\text{-}\acute{e}tale}}
\def\etale{{\acute{e}tale}}
\def\QCoh{\mathit{QCoh}}
\def\Ker{\mathop{\mathrm{Ker}}}
\def\Im{\mathop{\mathrm{Im}}}
\def\Coker{\mathop{\mathrm{Coker}}}
\def\Coim{\mathop{\mathrm{Coim}}}

% Boxtimes
%
\DeclareMathSymbol{\boxtimes}{\mathbin}{AMSa}{"02}

%
% Macros for moduli stacks/spaces
%
\def\QCohstack{\mathcal{QC}\!\mathit{oh}}
\def\Cohstack{\mathcal{C}\!\mathit{oh}}
\def\Spacesstack{\mathcal{S}\!\mathit{paces}}
\def\Quotfunctor{\mathrm{Quot}}
\def\Hilbfunctor{\mathrm{Hilb}}
\def\Curvesstack{\mathcal{C}\!\mathit{urves}}
\def\Polarizedstack{\mathcal{P}\!\mathit{olarized}}
\def\Complexesstack{\mathcal{C}\!\mathit{omplexes}}
% \Pic is the operator that assigns to X its picard group, usage \Pic(X)
% \Picardstack_{X/B} denotes the Picard stack of X over B
% \Picardfunctor_{X/B} denotes the Picard functor of X over B
\def\Pic{\mathop{\mathrm{Pic}}\nolimits}
\def\Picardstack{\mathcal{P}\!\mathit{ic}}
\def\Picardfunctor{\mathrm{Pic}}
\def\Deformationcategory{\mathcal{D}\!\mathit{ef}}
\usepackage{cleveref}
\crefname{originallemma}{original lemma}{original lemmas}
\Crefname{originallemma}{Original lemma}{Original lemmas}
\makeatletter
\newlabel{lemma-sheafification-tensor}{{0GMW}{source}{}{sourcepointer-1}{}}
\newlabel{lemma-sheafification-tensor@cref}{{[originallemma][0][]0GMW}{[1][0][]source}{}{}{}}
\makeatother
\title{Flatness of a sheafification from flatness of a cover of the presheaf of modules}
\author{The Stacks Project Authors}
\date{stacks-project:sites-modules.tex:local-snapshot-20261012; GNU FDL 1.2 or later}
\begin{document}
\maketitle
\noindent\textit{Editorial note.} All delivered mathematical text is the author's own, copied verbatim from the pinned Stacks Project chapter file sites-modules.tex; no statement or proof step was rewritten, reordered or removed.\par\medskip
\noindent\textit{Editorial note.} The counted span cites one related lemma, recorded in assembly.external\_references with its canonical tag as a navigation pointer.\par\medskip
\noindent\textit{Allowed standard prerequisites (editorial).}\begin{itemize}
\item Sheafification commutes with tensor products of presheaves of modules. Stacks Project, Modules on Sites, Lemma lemma-sheafification-tensor (tag 09BF). Identifies the tensor products of the sheafifications with sheafifications of presheaf tensor products, enabling the reduction.
\end{itemize}\medskip
\ifdefined\hypertarget\hypertarget{sourcepointer-1}{}\fi
\noindent\textit{Source reference.} Original lemma 0GMW: sources/stacks-chow/tags.txt, source lines 20771--20771. Cross-chapter standard prerequisite cited by the author inside the counted span; it is used as a pointer and its own proof is not part of this item.\par\medskip
\begin{lemma}
\label{lemma-flatness-sheafification-refined}
Let $\mathcal{C}$ be a site. Let $\mathcal{O}$ be a presheaf of rings.
Let $\mathcal{F}$ be a presheaf of $\mathcal{O}$-modules. Assume
that every object $U$ of $\mathcal{C}$ has a covering $\{U_i \to U\}_{i \in I}$
such that $\mathcal{F}(U_i)$ is a flat $\mathcal{O}(U_i)$-module.
Then $\mathcal{F}^\#$ is a flat $\mathcal{O}^\#$-module.
\end{lemma}

\begin{proof}
Let $\mathcal{G} \subset \mathcal{G}'$ be an inclusion of
$\mathcal{O}^\#$-modules. We have to show that
$$
\mathcal{G} \otimes_{\mathcal{O}^\#} \mathcal{F}^\#
\longrightarrow
\mathcal{G}' \otimes_{\mathcal{O}^\#} \mathcal{F}^\#
$$
is injective. By Lemma \ref{lemma-sheafification-tensor}
the source of this arrow is the sheafification of
the presheaf $\mathcal{G} \otimes_{p, \mathcal{O}} \mathcal{F}$
and similarly for the target. If $U$ is an object of $\mathcal{C}$ such that
$\mathcal{F}(U)$ is a flat $\mathcal{O}(U)$-module, then
$$
(\mathcal{G} \otimes_{p, \mathcal{O}} \mathcal{F})(U) =
\mathcal{G}(U) \otimes_{\mathcal{O}(U)} \mathcal{F}(U)
\longrightarrow
\mathcal{G}'(U) \otimes_{\mathcal{O}(U)} \mathcal{F}(U) =
(\mathcal{G}' \otimes_{p, \mathcal{O}} \mathcal{F})(U)
$$
is injective. Hence we reduce to showing: given a map of presheaves
$f : \mathcal{H} \to \mathcal{H}'$ on $\mathcal{C}$ such that
every $U$ in $\mathcal{C}$ has a covering $\{U_i \to U\}_{i \in I}$
with $\mathcal{H}(U_i) \to \mathcal{H}'(U_i)$ injective, then
$f^\#$ is injective. This we leave to the reader as an exercise.
\end{proof}

\end{document}
